\documentclass[11pt]{article}
\usepackage[margin=0.72in,top=0.62in,bottom=0.72in]{geometry}
\usepackage{graphicx}
\usepackage[dvipsnames]{xcolor}
\usepackage{comment}
\usepackage{csquotes}
\usepackage{amsmath}
\usepackage{fancyhdr}
\usepackage[normalem]{ulem}
\usepackage{bibentry}
\usepackage{enumitem}
\usepackage{titlesec}
\usepackage[colorlinks=true,urlcolor=LabOrange,linkcolor=LabTeal]{hyperref}

\definecolor{LabTeal}{HTML}{0C4050}
\definecolor{LabOrange}{HTML}{EA6738}
\definecolor{TextGray}{HTML}{26343A}
\definecolor{RuleGray}{HTML}{CBD5D8}
\definecolor{LabTint}{HTML}{F2F7F7}
\hypersetup{
  pdftitle={Keisuke Yoshihara - Academic Curriculum Vitae},
  pdfauthor={Keisuke Yoshihara},
  pdfsubject={Academic curriculum vitae}
}

\renewcommand{\familydefault}{\sfdefault}
\color{TextGray}
\setlength{\parindent}{0pt}
\setlength{\parskip}{1.5pt}
\setlist[description]{leftmargin=0.18in,labelsep=0.08in,itemsep=2pt,topsep=2pt}
\setlist[enumerate]{leftmargin=0.28in,itemsep=2pt,topsep=2pt}
\renewcommand{\labelenumi}{[\arabic{enumi}]}

\titleformat{\section}
  {\large\bfseries\color{LabTeal}}
  {}{0pt}{}
  [\vspace{-0.35em}\color{LabOrange}\rule{\textwidth}{1.2pt}]
\titlespacing*{\section}{0pt}{11pt}{5pt}
\titleformat{\subsection}{\normalsize\bfseries\color{LabTeal}}{}{0pt}{}

\pagestyle{fancy}
\fancyhf{}
\renewcommand{\headrulewidth}{0pt}
\renewcommand{\footrulewidth}{0.4pt}
\renewcommand{\footrule}{\hbox to\headwidth{\color{RuleGray}\leaders\hrule height \footrulewidth\hfill}\vskip 5pt}
\fancyfoot[L]{\footnotesize Keisuke Yoshihara | Academic Curriculum Vitae}
\fancyfoot[R]{\footnotesize \thepage}

\begin{document}
%\fancypagestyle{plain}

%%%%%%%%%%%%%%%%%%%%%
%%%% Title and Personal Info  %%%%%
%%%%%%%%%%%%%%%%%%%%%
\begin{minipage}[c]{0.78\textwidth}
{\fontsize{25}{28}\selectfont\bfseries\color{LabTeal} Keisuke Yoshihara}\\[3pt]
{\large\color{LabOrange} Experimental Particle Physicist}\\[7pt]
{\small University of Hawai'i at M$\bar{\text{a}}$noa}\\[2pt]
{\small +1 808 956 2941 \enspace | \enspace
\href{mailto:kyoshiha@hawaii.edu}{kyoshiha@hawaii.edu} \enspace | \enspace
\href{https://orcid.org/0000-0002-3656-2326}{ORCID 0000-0002-3656-2326}}
\end{minipage}
\hfill
\begin{minipage}[c]{0.16\textwidth}
\raggedleft\includegraphics[width=0.84in]{assets/yoshihara-group-logo.png}
\end{minipage}
\vspace{4pt}
\color{LabTeal}\rule{\textwidth}{2.2pt}
\color{TextGray}

\vspace{5pt}
\fcolorbox{RuleGray}{LabTint}{%
\parbox{\dimexpr\textwidth-2\fboxsep-2\fboxrule}{%
{\small\textbf{\color{LabTeal}Research profile:} Experimental particle physics at Belle II and SuperKEKB, spanning searches for physics beyond the Standard Model, detector operation and readout, FPGA firmware development, and open-source mixed-signal ASIC design. Current instrumentation work includes ultrafast beam diagnostics with the bunch oscillation recorder (BOR) and silicon-based X-ray radiation monitor (SiXRM).}%
}}

%%%%%%%%%%%%%%%%%%%%%
%%%%%%% Education   %%%%%%%%
%%%%%%%%%%%%%%%%%%%%%
%\vspace{-0.2in}
\section*{\large{Education}} 
{\footnotesize
\textbf{Doctor of Philosophy}, Physics, The University of Tokyo, 04/2011 - 09/2014 \\
Thesis: ``Measurement of the Higgs boson couplings using the $WW^{\ast}\rightarrow \ell\nu\ell\nu$ final state.'' 
[\htmladdnormallink{Link}{https://cds.cern.ch/record/1955933/}]\\ 
%supervised by Prof. T. Kawamoto \\
\textbf{Master of Science}, Physics, The University of Tokyo, 04/2009 - 03/2011,\\
Thesis: ``R$\&$D studies of Muon Polarimeter for J-PARC E06 T-Violation Experiment." 
[\htmladdnormallink{Link}{https://www.icepp.s.u-tokyo.ac.jp/papers/ps/thesis/master/kyoshihara_mthesis.pdf}]\\
%supervised by Prof. T. Kawamoto. \\
 \textbf{Bachelor of Science}, Physics, Yokohama National University, 04/2005 - 03/2009,\\
Thesis: ``Calculation of higher order diagrams for Compton Scattering on High Energy Limit." \\
%supervised by Prof. K. Sasaki. \\
}

%%%%%%%%%%%%%%%%%%%%%
%%% Professional Experience  %%%%%
%%%%%%%%%%%%%%%%%%%%%
\vspace{-0.2in}
\section*{\large{Professional Experience}}
{\footnotesize
\textbf{Assistant Professor}, University of Hawai'i at M$\bar{\text{a}}$noa, 04/2024 - present \\
\textbf{Visiting Researcher}, Nagoya University KMI, 05/2023 - present \\
\textbf{Data Science Consultant}, Accenture, 05/2023 - 03/2024 \\
\textbf{Visiting Associate Professor}, IPNS KEK, 04/2022 - 04/2023 \\
\textbf{Designated Associate Professor}, Nagoya University KMI, 04/2021 - 04/2023 \\
%\textbf{Postdoctoral Fellow}, IPNS KEK, 12/2019 - 03/2021 \\
\textbf{Postdoctoral Fellow}, Physics, Iowa State University, 10/2019 - 03/2021 \\
\textbf{Postdoctoral Fellow}, Physics, University of Pennsylvania, 10/2014 - 09/2019 
%\textbf{Research Assistant}, Physics, The University of Tokyo, 04/2011 - 09/2014 \\
%\textbf{Teaching Assistant}, Physics, KEK, 08/2009 KEK summer school \\ %; Course8: Superconducting RF cavity\\
}

%%%%%%%%%%%%%%%%%%%%%
%%%%%%% Awards etc %%%%%%%%
%%%%%%%%%%%%%%%%%%%%%
\section*{\large{Awards}} 

{\footnotesize
\begin{description}
\item[$\bullet$]{ATLAS Outstanding Achievement Award, ATLAS Collaboration, 01/2021
[\htmladdnormallink{Link}{https://atlas.cern/node/36504}]}
\item[$\bullet$]{Young Scientist Award of the Physical Society of Japan, JPS, 03/2016
[\htmladdnormallink{Link}{http://www.jps.or.jp/english/file/10th_wakate2016.pdf}]}
\item[$\bullet$]{Incentive Award of the Japan Association of High Energy Physicists, JAHEP, 10/2015
[\htmladdnormallink{Link}{http://www.jahep.org/office/syourei.html}]}
\item[$\bullet$]{Graduation with the highest honors (GPA 3.92/4.00), Yokohama National University, 03/2009}
\end{description}
}

\section*{\large{Research Grants}} 
{\footnotesize
\begin{description}
\item[$\bullet$]{Enhanced X-ray Beam Size Monitor (XRM) to Improve SuperKEKB Stability and Luminosity \\
\hspace{-0.25in} (Co-PI, U.S. DOE U.S.-Japan Science and Technology Cooperation Program in High Energy Physics, \$235,000) -- awarded, 2026}
\item[$\bullet$]{Construction of a Fast Beam Diagnosis System \\
\hspace{-0.25in} (PI, JSPS KAKENHI, 17.7 million JPY) -- awarded, 04/2022 [\htmladdnormallink{Link}{https://kaken.nii.ac.jp/ja/grant/KAKENHI-PROJECT-22H03867/}]} \\
\item[$\bullet$]{High Energy Physics Research at the University of Hawaii -- Feasibility Study of an Open-Source Mixed-Signal ASIC \\
\hspace{-0.25in} (Co-PI, DOE FOA-3432, \$75,000) -- awarded, 2026}
\end{description}
}

%%%%%%%%%%%%%%%%%%%%%
%%%%%%% Leadership %%%%%%%%
%%%%%%%%%%%%%%%%%%%%%
\vspace{-0.2in}
\section*{\large{Leadership}} 
{\footnotesize
\begin{description}
\item[$\bullet$]{Belle II Professional Development Committee Member (01/2026-present)}
\item[$\bullet$]{US Belle II DOE Level-2 Manager for TOP detector (07/2025-present)}
\item[$\bullet$]{US Belle II Executive Committee Member (01/2025-present)}
\item[$\bullet$]{Belle II Upgrade WG Coordinator of Outer Detector Sub-group, 09/2022 - 04/2023}
\item[$\bullet$]{Belle II (deputy) Run Coordinator, 05/2020 - 04/2023}
\item[$\bullet$]{Belle II KLM (deputy) Run Coordinator, 02/2020 - 03/2021}
\item[$\bullet$]{ATLAS SUSY Background Forum Convener, 10/2017 - 09/2019}
% -- convener takes responsibility on technical aspects of all analyses in ATLAS SUSY group
\item[$\bullet$]{ATLAS SUSY Stop1L Analysis Contact, 04/2016 - 10/2019} 
% -- analysis contact of the search for direct top squark pair production with 1-lepton final state 
\item[$\bullet$]{ATLAS TRT DAQ Operation Leader, 04/2017 - 03/2019}
 %-- convener (and team) is responsible for the entire hardware chain from the front-end to the back-end electronics
\item[$\bullet$]{ATLAS SUSY Higgsino Disappearing Track Analysis Contact, 10/2017 - 1/2018}
\item[$\bullet$]{ATLAS SUSY MC Generation Coordination, 11/2016 - 10/2017}
\item[$\bullet$]{ATLAS SUSY Fake-lepton Review Forum Coordinator, 05/2016 - 04/2017} \\
\end{description}
}

%%%%%%%%%%%%%%%%%%%%%
%%%%%%% Student supervision  %%%
%%%%%%%%%%%%%%%%%%%%%
\vspace{-0.2in}
\section*{\large{Mentoring at UH since 2024}} 
{\footnotesize
\begin{description}
\item[$\bullet$]{\textbf{Current Postdoctoral Researchers:} Tommy Lam, Horacio Crotte Ledesma, Martin Bessner, Shahab Kohani, Vasily Shebalin.}
\item[$\bullet$]{\textbf{Former Postdoctoral Researchers:} Christopher Ketter, Harsh Purwar.}
\item[$\bullet$]{\textbf{Graduate Students:} Cody Driver, Jacob Romeo, Robert Sinesi, Vikas Kumar.}
\item[$\bullet$]{\textbf{Undergraduate Students:} Aleena Bendzius, Yeunggyun Kwon.}
\end{description}
}

\begin{comment}
\vspace{-0.2in}
\section*{\large{Supervision for Ph.D. thesis}} 
{\footnotesize
For the following students, I have served as a mentor for more than 2 years to complete their Ph.D. thesis. I have also supervised many other graduate and undergraduate students in the area of hardware operation, development, and physics for short/long term.
\begin{description}
\item[$\bullet$] {P. Arrubarrena (LMU) graduated in 2020. Imperial College London postdoc:\\
\hspace{-0.25in} Machine learning in supersymmetry search}
\item[$\bullet$] {T. Yamazaki (UTokyo) graduated in 2020. UC Berkeley postdoc:\\
\hspace{-0.25in}Stop and higgsino search, development of soft b-tagging algorithm}
\item[$\bullet$] {D. Handle (LMU) graduated in 2019:\\
\hspace{-0.25in}Machine learning implementation in stop search}
\item[$\bullet$] {T. C. Herwig (UPenn) graduated in 2019. Fermilab postdoc:\\
\hspace{-0.25in}Study of ttbar/single-top Wt interference, search for stop decaying into higgsinos}
%\item {E. Reseguie (UPenn) graduated in 2019. LBNL Chamberlain Fellow:\\
%Search for electroweak SUSY production, TRT operation, ITK ASIC design}
%\item {M. Saito (UTokyo) graduated in 2019. UTokyo postdoc:\\
%Higgsino interpretation of disappearing track search}
\item[$\bullet$]{K. Onogi (Nagoya) graduated in 2018:\\
\hspace{-0.25in}Stop and higgsino searches in compressed scenarios, fake-lepton background study}
\item[$\bullet$]{J. Kuechler (Wuppertal) graduated in 2018. DESY Fellow:\\
\hspace{-0.25in}Stop search in boosted regime}
\item[$\bullet$]{D. Boerner (Wuppertal) graduated in 2018. DESY Fellow:\\
\hspace{-0.25in}Compressed stop search with MVA technique}
%\item{K. Brendlinger (UPenn) graduated in 2016. DESY postdoc:\\
%Electroweak SUSY search, fake-lepton background, electron identification efficiency}
%\item{R. Vanguri (UPenn) graduated in 2015:\\
%Search for Higgs boson decaying into invisible particles}
%\end{enumerate}
\end{description}
}
\end{comment}

%\noindent \underline{2. 研究グループリーダーとして指導・共同研究した博士課程学生 (名前と所属のみ)}\\
%2016: J. Gramling (Geneva), \\ %T. Mori (UTokyo), \\
%2017: A. Yiming (Stockholm), \\
%2018: A. Shcherbakova (Stockholm), \\
%2019: B. Haney (UPenn), K. Mistry (UPenn), J. Reichert (UPenn), J. Gonski (Harvard), \\
%2020 -- (not graduated yet): P. Arrubarrena (LMU), M. Claire (Cambridge), S. Andrean (Stockholm), L.~Sanchez (Stockholm), P.~Pasuwan (Stockholm) \\
%\vspace{-0.25in}

%%%%%%%%%%%%%%%%%%%%%
%%%%%%% Organization    %%%%%%
%%%%%%%%%%%%%%%%%%%%%
\section*{\large{Organization experience}} 
\hspace{-0.2in}
{\footnotesize
\begin{description}
\item[$\bullet$] {ML4FE Workshop 2025~[\htmladdnormallink{Link}{https://indico.phys.hawaii.edu/event/2506/}], 05/2025}
\item[$\bullet$] {2nd KMI Flash~[\htmladdnormallink{Link}{https://agenda.hepl.phys.nagoya-u.ac.jp/indico/conferenceDisplay.py?confId=2233}], 02/2023}
\item[$\bullet$] {Flavor Physics workshop 2022~[\htmladdnormallink{Link}{https://agenda.hepl.phys.nagoya-u.ac.jp/indico/conferenceDisplay.py?confId=2065}], \textbf{a main organizer}, 11/2022}
\item[$\bullet$] {1st KMI Flash~[\htmladdnormallink{Link}{https://agenda.hepl.phys.nagoya-u.ac.jp/indico/conferenceDisplay.py?ovw=True&confId=1774}], 07/2021}
\end{description}
I also have experience organizing workshops and meetings, and serving as a session chair at international conferences.
}

%%%%%%%%%%%%%%%%%%%%%
%%%%%%% Teaching %%%%%%%%%
%%%%%%%%%%%%%%%%%%%%%
\section*{\large{Teaching}} 
{\footnotesize
\begin{description}
\item[$\bullet$]  {``PHYS476: Modern Electronics for Physicists", UH Manoa, Spring 2025, Spring 2026. \\
\hspace{-0.25in}
Project-based course covering digital circuits, FPGA programming, Verilog HDL, and AI implementation on FPGA hardware.}
\item[$\bullet$]  {``PHYS399: Introduction to Scientific Programming", UH Manoa, Fall 2025. \\
\hspace{-0.25in}
Course introducing Python, C/C++, data analysis, large language model applications, and AI-agent development for scientific workflows.}
\end{description}
}

%%%%%%%%%%%%%%%%%%%%%
%%%%%%% Service %%%%%%%%
%%%%%%%%%%%%%%%%%%%%%
\begin{comment}
\section*{\large{Department Service}}
{\footnotesize
\begin{description}
\item[$\bullet$] {Graduate Program Advisory Committee Member, Department of Physics and Astronomy, UH Manoa (2024--present).}
\item[$\bullet$] {Contributed to graduate Core Competency Exam problem preparation, Department of Physics and Astronomy, UH Manoa (2025, 2026).}
\end{description}
}
\end{comment}

%%%%%%%%%%%%%%%%%%%%%
%%%%%%% Outreach   %%%%%%%%
%%%%%%%%%%%%%%%%%%%%%
\section*{\large{Outreach activities}} 
{\footnotesize
Selected outreach activities include:
\begin{description}
\item[$\bullet$]  {``Create a universe with a large accelerator - search for new physics beyond the Standard Model"[\htmladdnormallink{Link}{https://www.kmi.nagoya-u.ac.jp/wp-content/uploads/2021/07/A2poster-forPrint.pdf}],\\ 
\hspace{-0.25in}Science cafe jointly hosted by ICEHAP-KMI, 10/2021.} 
\item[$\bullet$]  {``Operation status of SuperKEKB/Belle II", KEK Press Conference, 08/2021.} 
\item[$\bullet$]  {``Approaching mysteries of the universe with accelerator", The 1st KMI-ITbM Mix Cafe, 05/2021.} 
\item[$\bullet$]  {``Introduction to particle physics", Shibuya Makuhari high school, 10/2020.} 
\item[$\bullet$]  {``Come on B-mesons!"[\htmladdnormallink{Link}{https://www.youtube.com/watch?v=SEJ6lKP41kY&t=1107s}], KEK open house, 09/2020.} 
\item[$\bullet$]  {``Let's do science with Haiku", NHK E-television, 08/2020.} 
\end{description}
}

%%%%%%%%%%%%%%%%%%%%%
%%%%%%% Presentation  %%%%%%%
%%%%%%%%%%%%%%%%%%%%%
%\newpage
\section*{\large{Selected Presentations since 2024}} 
%\vspace{0.2in}
%\hspace{0.2in}
%{\textbf{Physics Seminar/Colloquium and Invited Talks}}
%{\footnotesize
%\begin{enumerate}
%\renewcommand{\labelenumi}{[\arabic{enumi}]} 
%\item {``Current Status and Prospects of B-Physics",\\ \textbf{KMI Symposium 2025}, Nagoya, Japan, 03/2025.} 
%\item {``Challenges towards the Luminosity Frontier at the SuperKEKB/Belle II Experiment", \\ \textbf{University of Hawaii at Manoa}, HI, US, 04/2023.} 
%\item {``Current Status of New Physics Beyond the Standard Model in the Belle II Experiment and Prospects After the SuperKEKB Upgrade", \\ \textbf{JPS symposium(ONLINE)}, Japan, 03/2023.} 
%\item {``B-physics at Belle II experiment", \\ \textbf{KMI Symposium 2023}, Nagoya, Japan, 02/2023.} 
%\item {``Status and prospects for SuperKEKB/Belle II Experiment", \\ \textbf{University of Zurich (ONLINE)}, Zurich, Switzerland, 09/2022.} 
%\item {``SuperKEKB/Belle II Experiment: Status and Prospect", \\ \textbf{Iowa State University}, IA, US, 08/2022.} 
%\item {``Search for new physics at intensity frontier", \\ \textbf{JPS symposium (ONLINE)}, Japan, 03/2022.} 
%\item {``Status of the Belle II experiment", \\ \textbf{ILC summer camp 2021 (ONLINE)}, Japan, 08/2021.} \item {``Status and prospect for SuperKEKB/Belle II Experiment", \\ %\textbf{Nagoya University, KMI Topics seminar (ONLINE)}, Japan, 07/2021.} 
%\item {``Status of Belle II/SuperKEKB operation", \\ \textbf{23rd meeting on physics at B factories (ONLINE)}, Japan, 12/2020.} 
%\item {``Cornering SUSY naturalness with the ATLAS experiment", \\ \textbf{Princeton University}, NJ, US, 10/2017.} 
%\item {``Search for 3rd gen. squarks and cornering natural SUSY with the ATLAS detector", \\ \textbf{Fermilab}, IL, US, 10/2017, \textbf{(Fermilab Wine and Cheese Seminar)}.}
%\item {``Cornering SUSY naturalness with the ATLAS experiment", \\ \textbf{Johns Hopkins University}, MD, US, 10/2017.} 
%\item {``Cornering SUSY naturalness with the ATLAS experiment", \\ \textbf{University of Pittsburgh}, PA, US, 09/2017.} 
%\item {``Going beyond the Higgs discovery at the LHC", \\ \textbf{Osaka University}, Japan, 07/2017.} 
%\item {``Recent results on SUSY searches in LHC-ATLAS", \\ \textbf{University of Pennsylvania}, PA, US, 04/2016.}
%\item {``Recent results on SUSY searches from LHC-ATLAS", \\ \textbf{KEK}, Tsukuba, Japan, 03/2016.}
%\item {``Measurement of the Higgs boson couplings using the $WW^{\ast}\rightarrow \ell\nu\ell\nu$ final state.", \\ \textbf{JPS}, Sendai, Japan, 03/2016, \textbf{(Invited as a winner of Young Scientists Award)}.}
%\item {``Evidence for the Higgs Boson in the H to $WW^{\ast} \rightarrow l\nu l\nu$ Channel from the ATLAS Experiment at the LHC", \textbf{University of Pennsylvania}, PA, US, 11/2013.}
%\item {``Evidence for the Higgs boson and Measurement of its properties in the $WW^{\ast} \rightarrow l\nu l\nu$ decay with the ATLAS", \textbf{University of Illinois at Chicago}, IL, US, 11/2013.}\\
%\end{enumerate}
%}

%International Conference
%{\textbf{Presentations at conferences and workshops since 2024}}
{\footnotesize
\begin{enumerate}
\renewcommand{\labelenumi}{[\arabic{enumi}]} 
\item {``AI/ML on Hardware'',\\
\textbf{Belle II Summer Workshop}, University of Louisville, KY, US, 07/2026.}
\item {``Machine Detector Interface'',\\
\textbf{Belle II Summer Workshop}, University of Louisville, KY, US, 07/2026.}
\item {``The Belle II Upgrade Program'',\\
\textbf{Violation of Fundamental Symmetries with B mesons Workshop}, Fermilab, IL, US, 2025.}
\item {``AI/ML-Based Beam Diagnostics'',\\
\textbf{CEF Workshop: Versal and AI in FEE}, KEK, Japan, 2025.}
\item {``Front-end AI/ML Opportunities at SuperKEKB/Belle II'',\\
\textbf{CEF Workshop: Versal and AI in FEE}, KEK, Japan, 2025.}
\item {``Modern Electronics Education with AI/ML and FPGA'',\\
\textbf{CPAD2025 Workshop}, University of Pennsylvania, PA, US, 2025.}
\item {``Introduction to SuperKEKB/Belle~II Machine Detector Interface'',\\
\textbf{Belle II Summer Workshop}, Virginia Tech, VA, US, 2025.}
\item {``Modern Electronics for Physicists: Project-Based Learning in PHYS476",\\ \textbf{US-Japan Symposium}, Sapporo, Japan, 05/2025.} 
\item {``Development and Implementation of Advanced Beam Diagnostic and Abort Systems in SuperKEKB",\\ \textbf{US-Japan Symposium}, Sapporo, Japan, 05/2025.} 
\item {``Modern Electronics for Physicists: Project-Based Learning in PHYS476",\\ \textbf{ML4FE workshop}, HI, US, 05/2025.} 
\item {``Current Status and Prospects of B-Physics",\\ \textbf{KMI Symposium 2025}, Nagoya, Japan, 03/2025.} 
\item {``Belle~II Loss Monitor Group's Activities for Understanding the SBL Mechanism",\\ \textbf{eeFACT conference}, Tsukuba, Japan, 03/2025.} 
\item {``The Belle~II Upgrade Program",\\ \textbf{CPAD workshop 2024}, TN, US, 11/2024.} 
%\item{``Status and prospects for SuperKEKB/Belle II Experiment", \\ \textbf{APPC15}, Seoul, Korea.}\label{APPC15} 
%\item{``Belle II: Analysis Summary of $b \rightarrow s$, $b \rightarrow s+\gamma$, and $b \rightarrow s\ell^+\ell^-$ Transitions", \\ \textbf{BEAUTY2020}, Chiba, Japan.}\label{BEAUTY2020} 
%\item{``Search for direct stop pair production with the ATLAS detector", \\ \textbf{DPF2019}, MA, US.}\label{DPF2019} 
%\item{``Search for direct pair production of 3rd generation squarks with the ATLAS detector", \\ \textbf{EPS2019}, Ghent, Belgium.}\label{EPS2019} 
%\item{``Search for 3rd generation superpartners with the ATLAS experiment", \\ \textbf{DPF2017}, IL, US.}\label{DPF2017} 
%\item{``Searches for third generation squarks with leptonic channels with the ATLAS", \\ \textbf{LHCP2017}, Shanghai, China.}\label{LHCP2017}
%\item{``SM Higgs Combination and Higgs Properties Measurements in ATLAS", \\ \textbf{SUSY 2014}, Manchester, UK.}\label{SUSY2014} 
%\item{``Measurements of the properties of the Higgs boson in the $H\rightarrow WW^{\ast}\rightarrow l\nu l\nu$ decay channel with the ATLAS detector using 25 $\mathrm{fb^{-1}}$ of proton-proton collision data", \\ \textbf{Lepton Photon 2013}, San Francisco, US, Poster Session.}\label{LP2013} 
%\item{``Update of the $H\rightarrow WW^{\ast}\rightarrow e\nu \mu\nu$ analysis with 13 $\mathrm{fb^{-1}}$ of $\sqrt{s} = 8$ TeV data collected with the ATLAS detector", \textbf{Hadron Collider Physics Symposium 2012}, Kyoto, Japan, Poster Session.}%\label{HCP2012} \\
\end{enumerate}
}

%%%%%%%%%%%%%%%%%%%%%
%%%%%%% Publication %%%%%%%%
%%%%%%%%%%%%%%%%%%%%%
\section*{\Large{Publications since 2024}} 
{\normalsize\bfseries\color{LabOrange}Peer-reviewed Journals and Preprints}
\vspace{0.04in}
{\footnotesize
\begin{enumerate}
\renewcommand{\labelenumi}{[\arabic{enumi}]} 
\item{H. Purwar \textit{et al.}, ``Migration of Belle II TOP Feature Extraction from the Zynq Processing System to PCIe40 Readout PCs'', arXiv:2608.27814 (2026), submitted to \textit{Nucl. Instrum. Methods Phys. Res. A}. \href{https://arxiv.org/abs/2608.27814}{arXiv}}
\item{Y.-T. Lai, T. Koga, Y. Nakazawa, N. Taniguchi, and \underline{K. Yoshihara}, ``Machine-Learning-Based Waveform Discrimination in the Front-End Electronics of the Belle II Central Drift Chamber for Cross-Talk Noise Reduction'', arXiv:2607.25174 (2026), submitted to \textit{IEEE Trans. Nucl. Sci.}. \href{https://arxiv.org/abs/2607.25174}{arXiv}}
\item{M. Andrew \textit{et al.}, ``Development of a Silicon-Based Ultra-Fast X-Ray Beam Size Monitor for SuperKEKB'', arXiv:2607.08205 (2026), submitted to \textit{Nucl. Instrum. Methods Phys. Res. A}. \href{https://arxiv.org/abs/2607.08205}{arXiv}}
\item{J. Gonski \textit{et al.}, ``Machine Learning on Heterogeneous, Edge, and Quantum Hardware for Particle Physics'', \textit{PRX Intelligence}, accepted 31 July 2026. \href{https://doi.org/10.1103/bjx3-vlj8}{DOI}}
\item{C. Ketter \textit{et al.}, ``Design and Commissioning of Readout Electronics for a $K_L^0$ and $\mu$ Detector at the Belle II Experiment'', \textit{Nucl. Instrum. Methods Phys. Res. A} \textbf{1082} (2026) 170893. \href{https://doi.org/10.1016/j.nima.2025.170893}{DOI}}
\item{H. Atmacan \textit{et al.}, ``The Imaging Time-of-Propagation Detector at Belle II'', \textit{Nucl. Instrum. Methods Phys. Res. A} \textbf{1080} (2025) 170627. \href{https://doi.org/10.1016/j.nima.2025.170627}{DOI}}
\item{\underline{K. Yoshihara} \textit{et al.}, ``Development and Implementation of Advanced Beam Diagnostic and Abort Systems in SuperKEKB'', \textit{Nucl. Instrum. Methods Phys. Res. A} \textbf{1072} (2025) 170117. \href{https://doi.org/10.1016/j.nima.2024.170117}{DOI}}
\end{enumerate}

{\normalsize\bfseries\color{LabOrange}Conference Proceedings}
\begin{enumerate}
\renewcommand{\labelenumi}{[\arabic{enumi}]}
\item{R. Nomaru \textit{et al.}, ``Development and First Beam Observations of an Ultra-Fast Bunch-by-Bunch X-Ray Beam Size Monitor at SuperKEKB'', in \textit{Proc. IPAC'26}, Deauville, France, 2026, paper MOV6603. \href{https://doi.org/10.18429/JACoW-IPAC2026-MOV6603}{DOI}}
\item{R. Nomaru \textit{et al.}, ``R\&D of an Ultrafast X-Ray Beam Size Monitor for SuperKEKB'', in \textit{Proc. IBIC'25}, Liverpool, UK, 2025, pp. 435--438. \href{https://doi.org/10.18429/JACoW-IBIC2025-TUPCO31}{DOI}}
\item{R. Zhang \textit{et al.}, ``Study on a Novel Laser Fast Abort System for SuperKEKB'', in \textit{Proc. IPAC'25}, Taipei, Taiwan, 2025, pp. 3048--3050. \href{https://doi.org/10.18429/JACoW-IPAC2025-THPS038}{DOI}}
\item{H. Kaji \textit{et al.}, ``Upgrade of Beam Abort System at the SuperKEKB Positron Ring'', in \textit{Proc. IPAC'25}, Taipei, Taiwan, 2025, pp. 3037--3039. \href{https://doi.org/10.18429/JACoW-IPAC2025-THPS033}{DOI}}
\item{\underline{K. Yoshihara} \textit{et al.}, ``Loss Monitor Group's Activities for Understanding the SBL Mechanism at SuperKEKB/Belle II'', in \textit{Proc. eeFACT'25}, Tsukuba, Japan, 2025.}
\end{enumerate}
}

\begin{comment}

%Regional Conferences
{\textbf{Regional conferences}}
{\footnotesize
\begin{enumerate}
\renewcommand{\labelenumi}{[\arabic{enumi}]} 
\item{\underline{K. Yoshihara} for SuperKEKB/Belle~II Collaboration, ``Development of beam diagnostic system for the SuperKEKB accelerator", Japanese Physical Society Meeting (JPS), ONLINE, Japan, 03/2022.}
\item{\underline{K. Yoshihara} for Belle Collaboration, ``Search for $B$-meson rare decay in $b \rightarrow d \ell^{+}\ell^{-}$ channel with the Belle detector", Japanese Physical Society Meeting (JPS), Nagoya, Japan, 03/2020.}
\item{\underline{K. Yoshihara} for ATLAS Collaboration, ``Measurement of the Higgs couplings in the $H\rightarrow WW^{(\ast)}\rightarrow l\nu l\nu$ decay channel with the ATLAS detector", JPS, Sagamihara, Japan, 03/2014.}
\item{\underline{K. Yoshihara} for ATLAS Collaboration, ``Measurement of the Higgs couplings in the $H\rightarrow WW^{(\ast)}\rightarrow l\nu l\nu$ decay channel with the ATLAS detector", JPS, Hiroshima, Japan, 03/2013. Similar talks given in the JPS, 09/2012, 03/2012 and 09/2011.}
\item{\underline{K. Yoshihara} for ATLAS Collaboration, ``Search for the Higgs boson in the $H\rightarrow WW^{(\ast)}\rightarrow l\nu l\nu$ decay channel with the ATLAS detector", JPS, Kyoto, Japan, 09/2012.}
\item{\underline{K. Yoshihara} for ATLAS Collaboration, ``Search for the Higgs boson in the $H\rightarrow WW^{(\ast)}\rightarrow l\nu l\nu$ decay channel with the ATLAS detector", JPS, Hirosaki, Japan, 03/2012.}
\item{\underline{K. Yoshihara} for ATLAS Collaboration, ``Search for the Higgs boson in the $H\rightarrow WW^{(\ast)}\rightarrow l\nu l\nu$ decay channel with the ATLAS detector", JPS, Nishinomiya, Japan, 09/2011.}
\item{\underline{K. Yoshihara}, A. Toyoda, F. Habermehl et al., ``First Commissioning of the K1.1BR Beamline at Hadron Hall in J-PARC", JPS, Nigata, Japan, 03/2011.}
\item{\underline{K. Yoshihara}, Y. Igarashi, J. Imazato et al., ``Study of Muon Polarimeter for J-PARC E06 Experiment using muon beam at TRIUMF", JPS, Kitakyushu, Japan, 09/2010. Similar talk given in the JPS, 03/2010.} 
\item{\underline{K. Yoshihara}, Y. Igarashi, J. Imazato et al., ``Study of Muon Polarimeter for J-PARC E06 Experiment using muon beam at TRIUMF", JPS, Okayama, Japan, 03/2010.}
\end{enumerate}
}

\newpage

%%%%%%%%%%%%%%%%%%%%%
%%%%     Research Experience   %%%%
%%%%%%%%%%%%%%%%%%%%%
%\hrule height 0.5mm depth 0.5mm width 160mm
\section*{\large{\uline{Research Experience}}} 

%%%%%%%%%%%%%%%%%%%%%
%%%%%%% Belle/Belle II  %%%%%%%
%%%%%%%%%%%%%%%%%%%%%
{\small

\subsection*{\large \bf{1. SuperKEKB/Belle II experiment}}

\paragraph{KLM operation} 02/2020 - 03/2021 \\
I was involved in the operation of ``$K_{L}^{0}$ and $\mu$" (KLM) detector and stepped up as Deputy KLM run coordinator in February 2020, taking responsibilities on (1) diagnosing and solving issues during the operation, (2) coordinating local activities, and (3) acting as a contact person to members of the KLM group and to those in other sub-systems. During the early 2020 operation, frequent errors (ttlost and b2llost) of the readout system have occurred posing the Belle II DAQ. I have led the effort in identifying the problem, which was turned out to be attributed to the front-end communication, and in modifying the configuration of the readout scheme by (1) shortening the cable length, (2) electrical noise reduction by a better grounding of the VME backplane, and (3) manufacturing of 4-slot VME crate that can fit in a narrow space in the detector area. In the late 2020 operation, the error rate has been significantly reduced (90\% reduction) after the modification and the stable operation and data-taking were achieved. I have also developed an online monitoring system with Elasticsearch which helped find problematic hardware channel immediately.

\paragraph{Belle II run coordination} 05/2020 - 04/2023 \\
I have been appointed as Deputy Run Coordinator in May 2020 and have taken a lead in SuperKEKB/Belle II operation in collaboration with K. Matsuoka (KEK) and H. Atmacan (Cincinnati). The Deputy Run Coordinator is tasked with taking a lead in (1) communication with accelerator group in daily basis to make a daily operation plan and in (2) diagnosis of detector issues working closely with Belle II sub-system experts. In the operation, there are two different operation shift -- Control Room shift and Belle II Commissioning Group (BCG) shift. Control Room shifters and BCG shifter work together sharing responsibilities on the Belle II operation. BCG shifter is particularly responsible for online monitoring of beam background and is served as a commander for safe operation of the Belle II HV system. I have prepared a shift manual for BCG and support shifters online in 24/7. 

\paragraph{Beam background mitigation with movable collimator} 03/2020 - 06/2022 \\
I have taken responsibility on the operation of vacuum collimators with S. Tanaka (KEK) and I. Nakamura (KEK). The collimators play an essential role to protect the Belle II detector from the beam background and their aperture is an important parameters for accelerator tuning. Non optimal setting of the aperture may result in 5-10 times higher level-1 trigger rate due to the beam background. Given that the machine condition and the beam optics keep changing during the operation, I have made sure the beam background being constant while retaining the injection efficiency and the lifetime of the beam. 

\paragraph{Development of beam diagnosis system} 06/2021 - present \\
As the beam current increases, a significant beam loss accompanying with quenching of the final focusing magnet has been observed due to the instability of the beam which was caused by unknown reasons. Therefore the increase of the beam current has been very challenging in the machine operation. In order to identify the location of the initial beam loss, a new group has been formed jointly with the accelerator group and I have led the group putting an effort for the development of the beam diagnosis system. The system consists of fast loss monitors and high precision timing synchronization system called White Rabbit~[\htmladdnormallink{Link}{https://white-rabbit.web.cern.ch/}], which is synchronized with the GPS, providing the time stamp of the loss monitors being fired. We were finally able to determine the position of the initial loss with the system.


%%%%%%%%%%%%%%%%%%%%%
%%%%%%% ATLAS %%%%%%%%%%
%%%%%%%%%%%%%%%%%%%%%
{\small

\subsection*{\large \bf{2. LHC-ATLAS experiment}}
{\small

\paragraph{Supersymmetry searches} 10/2014 - 04/2020 \\
%Since the discovery of the SM Higgs boson, the hierarchy problem has been a significant focus, highlighting the divergent corrections to the Higgs mass at high energy scales. SUSY, a promising theoretical framework, addresses this issue along with other unresolved problems in the SM. The superpartner of the top quark (stop) plays a crucial role in canceling the mass divergence due to the top quark's large Yukawa coupling. 
Many natural SUSY models predict a light stop with a mass under 1 TeV. From October 2014 to the April 2020, I have been deeply involved in the search for the top quark's superpartner, co-coordinating the stop1L analysis team since 2016. This team, consisting of over thirty physicists, searches for direct stop pair production in events with one prompt lepton, jets, and significant missing transverse energy, covering a wide range of potential stop and LSP masses. I introduced new signal models and developed a novel stop-higgsino analysis with students, utilizing low-$p_T$ leptons. Additionally, we employed machine learning techniques, specifically Recurrent Neural Networks (RNNs), to enhance sensitivity in discriminating stop signals from $t\bar{t}$ events. This analysis, based on Keras and TensorFlow, was the first in ATLAS to utilize ML algorithms. In scenarios with very small mass-splitting between the stop and LSP, we developed a secondary vertex reconstruction algorithm for soft b-tagging, significantly lowering the $p_T$ threshold for b-hadron reconstruction. Furthermore, I extended the disappearing track analysis to probe higgsinos with small mass-splitting, developing a higgsino model and co-coordinating the analysis team. My role as a co-convener of the SUSY Background Forum and SUSY MC production coordinator involved overseeing common physics objects and methods for background determination, as well as generating numerous signal and background samples. My expertise in fake-lepton backgrounds from the Higgs to $WW^{*}$ analysis led to co-leading the SUSY fake-lepton review forum, developing methods and tools recognized across the collaboration.

\paragraph{Higgs to  $WW^{*}$ analysis} 05/2011 - 09/2014 \\
I was an active member of the Higgs to $WW^{*} \rightarrow l \nu l \nu$ analysis group until 2014. One of the significant backgrounds in this analysis is the $W$+jets process, where $W$ bosons are produced with a jet that is misidentified as a lepton, known as a fake-lepton. Although not the largest background, the systematic uncertainty on the fake-lepton background was substantial (around 50\%), limiting experimental sensitivity. From May 2011 to July 2012, I collaborated with the ATLAS e/gamma group on measuring the efficiency of electron identification using $Z \rightarrow ee$ events. This measurement was critical for the Higgs boson discovery in both the $WW$ and $ZZ$ channels, but low-$p_T$ measurements were hindered by background contamination. I addressed this by introducing a data-driven template method with same-charge dilepton events, providing a robust cross-check to the baseline method. From May 2011 to September 2014, I developed a data-driven fake-factor method for the fake-lepton background, working with the Penn group. The dominant uncertainty came from extrapolating the quark-gluon composition from the control region to the signal region. By analyzing the origin of fake-leptons and their charge correlation in same-charge dilepton events, I reduced the systematic uncertainty on this extrapolation, allowing for a lower subleading lepton $p_T$. This improvement in reducing systematic uncertainty and lowering the lepton $p_T$ threshold increased the signal significance by 18\%. Combined with other advancements, this led to the world's best measurement of Higgs boson couplings in a single channel at the LHC during Run-1. I received my PhD based on this analysis, as well as the Young Scientist Award from the Japanese Physical Society.

\paragraph{TRT detector operation} 02/2016 - 01/2019 \\
The ATLAS TRT is the outermost inner detector subsystem, providing additional discrimination against hadrons in the identification of electrons. Since 2017, I have been co-coordinating the operation of the ATLAS TRT detector, first with B.~Haney (Penn), second with E.~Resseguie (Penn), and now with S. Chen (Penn), taking the responsibility for the full TRT readout chain, both electronics and the TRT DAQ system. %~[\htmladdnormallink{Link}{https://cds.cern.ch/record/2286227/files/ATL-INDET-INT-2017-015.pdf}].} 
Our responsibilities include diagnosing problems as they arise, replacing failed off-detector electronics (e.g. Readout Driver Boards (ROD's), Timing and Trigger Control Boards (TTC's), power supply cards, or receiver-driver boards at the PP2 location on the detector) and maintaining, via repair and testing, an adequate supply of spares for the PP2 receiver-driver cards. We have also been in charge of the calibration of the TRT straw hits, firmware testing, and maintenance of the TRT software to accommodate changes in the ATLAS central DAQ software. This work is essential to smooth ATLAS data-taking. Operational efficiencies near 100\% have been achieved so far, despite the challenges of increasing luminosity and the high-pileup environment.

\paragraph{Design of ASICs for ATLAS Inner Tracker} 10/2014 - 09/2015 \\
The HL-LHC is a future project of the LHC, where the operation of the ATLAS detector will be extremely challenging because of the increased luminosity. Therefore, it is important to upgrade a large part of the detector, and in particular to develop a new all-silicon Inner Tracker (ITK), composed of pixel and strip layers, a centerpiece of the ongoing upgrade activities for ATLAS. The Penn group has been significantly involved in the ITK strip ASIC group, working on the design of three ASIC chips [\htmladdnormallink{Ref}{https://cds.cern.ch/record/2257755/files/ATLAS-TDR-025.pdf}]: the Analogue binary chip (ABCStar), Hybrid control chip (HCCStar), and Autonomous Monitoring and Control (AMAC) chip. I have studied the behavioral modeling of the chips in hardware simulation, implementing all functionalities (logic circuits) and optimizing the data packet structure for the ABC and HCC to accommodate trigger-rate and latency constraints. I have found that having a variable length structure in the HCC improves the speed of data transfer significantly, and that the strip readout will be able to accommodate a high trigger rate (up to L0 = 1~MHz) without major technical problems. However, given the complexity of the handling of variable length packets, we decided to stick to the fixed length packet structure.
}

%%%%%%%%%%%%%%%%%%%%%
%%%%%%% J-PARC   %%%%%%%%%
%%%%%%%%%%%%%%%%%%%%%
{\small

\subsection*{3. Kaon Experiments at J-PARC (2009-2011) }
I have worked on two kaon experiments at J-PARC: the E36 experiment [\htmladdnormallink{Link}{http://j-parc.jp/researcher/Hadron/en/pac_1007/pdf/KEK_J-PARC-PAC2010-04.pdf}] aimed for a precision measurement of lepton universality, and the E06 experiment [\htmladdnormallink{Link}{http://trek.kek.jp/proposal/E06_proposal.pdf}] targets a direct T-violation search in muon transverse polarization ($P_{\mathrm{T}}$) in $K^+ \rightarrow \pi^0 \mu^+ \nu_{\mu}$ decays. 
%Most of my time was devoted for the R\&D studies of detectors necessary for these experiments and a beam commissioning at J-PARC.

\paragraph{Development of muon polarimeter} 04/2009 - 03/2011 \\
I developed a novel muon polarimeter in collaboration with J.~Imazato (KEK) and the REPIC company. The polarimeter is a drift tube tracking chamber, precisely measuring the direction of the positron emitted from the muon decay. The principle of the polarimeter is based on the fact that the positron preserves the muon spin information due to helicity conservation. The polarimeter needs some amount of material to stop incoming muons in contrast to normal drift chambers where the amount of the inside material is minimized to suppress multiple scattering effects. The inner structure of the polarimeter was optimized using GEANT4 simulation. The material of the tubes was selected carefully, taking into account the effect of muon spin rotation ($\mu$SR) arising from external magnetic fields produced by nuclei in the material, as this potentially changes the muon spin direction. 

\paragraph{Commissioning of K1.1BR beamline at J-PARC} 07/2010 - 03/2011 \\
For the K1.1BR beamline commissioning, I prepared and installed various detectors nessessary for the beam monitoring in collaboration with Y.~Igarashi (KEK) and J.~Imazato (KEK). These include (1) a Hodoscope counter for the beam monitoring, (2) a Fitch-type Cherenkov counter essential to separate kaons from pions, which make up an extremely large fraction the beam, (3) Time-of-Flight counters for $\beta$ measurement, (4) MWPCs for the beam profiling, and (5) an Air Cherenkov counter for positron discrimination. The Hodoscope was the first detector that I fully designed and constructed myself. 
%It consists of $\sim$ 24 plastic scintillation counters and 1/2 inch PMTs. 
I also worked on the DAQ system and provided various triggers (proton, kaon, pion, muon and positron) for the beamline users for test experiments. I also contributed to the optimization study of the dipole and quadrupole magnets in the beamline. The extraction of a high rate/purity kaon beam was achieved. The R$\&$D studies mentioned above are summarized in my Master of Science thesis [\htmladdnormallink{Ref}{http://www.icepp.s.u-tokyo.ac.jp/papers/ps/thesis/master/kyoshihara_mthesis.pdf}]. 

%\begin{itemize}
%\item {I have contributed to the R$\&$D studies of CsI(Tl) calorimeter readout chips such as Avalanche Photo Diodes (APD) and Micro pixel Avalanche Photo Diodes (MAPD). I led the effort in preparing the experimental setup for the beam test, including data acquisition systems.}
%\end{itemize}
}

\end{comment}

\end{document}
